\part{Other principal \(q\)-analogs}

\section{\(q\)-numbers}
\label{sec:qs-qnumbers}

Choose a logarithm value \(\PrincipalLogarithm q\) for \(q\in\ComplexNumbers\setminus\{0,1\}\), put
\(q^x=\EulerE^{x\PrincipalLogarithm q}\), and define
\[
  [x]_q:=\frac{1-q^x}{1-q}.
\]

\begin{proposition}[Exact inverse in the exponent]
\label{prop:qs-qnumber-x}
Let \(1-(1-q)y\ne0\), and choose a value of
\(\PrincipalLogarithm(1-(1-q)y)\).  Then
\[
  x=\frac{\PrincipalLogarithm(1-(1-q)y)}{\PrincipalLogarithm q}
\]
solves \([x]_q=y\).  Every solution is obtained by allowing the logarithm
value in the numerator to vary by \(2\pi\ImaginaryUnit\IntegerNumbers\).
\end{proposition}

\begin{proof}
The equation \([x]_q=y\) is equivalent to
\[
  \EulerE^{x\PrincipalLogarithm q}=1-(1-q)y.
\]
Taking a logarithm value gives the displayed expression.  Conversely,
exponentiating that expression recovers the equation.  The complete
preimage of a nonzero complex number under the exponential consists of one
logarithm value plus \(2\pi\ImaginaryUnit\IntegerNumbers\), which proves the last assertion.
\end{proof}

For an integer \(m\ge2\),
\[
  [m]_q=1+q+\cdots+q^{m-1},
\]
so inversion in the base is an algebraic covering.

\begin{proposition}[Distinguished base jets of the \(q\)-number]
\label{prop:qs-qnumber-jets}
Let \(u=y-1\).  The inverse branch through \(q=0\) satisfies
\[
\begin{array}{rll}
 m=2:&q=u,&\text{exactly},\\
 m=3:&q=u-u^2+2u^3+\BigO(u^4),&\\
 m\ge4:&q=u-u^2+u^3+\BigO(u^4).&
\end{array}
\]
For a fixed complex \(x\) and \(q=\EulerE^{-t}\to1\),
\[
  [x]_{\EulerE^{-t}}
  =x+\frac{x(1-x)}2t
   +\frac{x(2x^2-3x+1)}{12}t^2
   -\frac{x^2(x-1)^2}{24}t^3+\BigO(t^4).
\]
If \(x\notin\{0,1\}\), its base inverse is
\[
  t=\frac{2(y-x)}{x(1-x)}
   +\frac{2(2x-1)}{3x^2(x-1)^2}(y-x)^2
   +\BigO((y-x)^3).
\]
\end{proposition}

\begin{proof}
At zero, the forward series begins \(u=q\) for \(m=2\),
\(u=q+q^2\) for \(m=3\), and \(u=q+q^2+q^3+\BigO(q^4)\) for
\(m\ge4\).  Direct substitution of
\(q=u+Au^2+Bu^3+\cdots\) gives the three formulas.

For \(q=\EulerE^{-t}\), divide the Taylor series of
\(1-\EulerE^{-xt}\) by that of \(1-\EulerE^{-t}\), or equivalently use
\[
  [x]_{\EulerE^{-t}}
  =x\exp\left(
    -\frac{x-1}{2}t+
    \sum_{r\ge1}
    \frac{B_{2r}(x^{2r}-1)}{2r(2r)!}t^{2r}
  \right).
\]
Expansion through degree three gives the stated forward jet.  Reversion of
\(\delta:=y-x=a t+b t^2+\cdots\), where
\[
  a=\frac{x(1-x)}2,\qquad
  b=\frac{x(x-1)(2x-1)}{12},
\]
gives \(t=\delta/a-b\delta^2/a^3+\cdots\), which is the displayed inverse.
For \(x=0\) or \(1\), the forward function is identically \(0\) or \(1\), so
base recovery is impossible rather than singular of finite order.
\end{proof}

\begin{proposition}[Alternating and reciprocal \(q\)-number branches]
\label{prop:qs-qnumber-alternating}
For \(r\ge1\),
\[
  [2r]_{-1}=0,\qquad
  \left.\frac{\partial}{\partial q}[2r]_q\right|_{q=-1}=r,
\]
and
\[
  [2r+1]_{-1}=1,\qquad
  \left.\frac{\partial}{\partial q}[2r+1]_q\right|_{q=-1}=-r.
\]
At infinity,
\[
  [m]_q=q^{m-1}(1+q^{-1}+\cdots+q^{-(m-1)}),
\]
so the inverse at a large target has \(m-1\) Puiseux determinations, indexed
by the choices of \(y^{1/(m-1)}\).
\end{proposition}

\begin{proof}
Evaluate \(\sum_{j=0}^{m-1}q^j\) and its derivative
\(\sum_{j=1}^{m-1}jq^{j-1}\) at \(-1\), pairing consecutive terms.
The reciprocal identity is obtained by factoring out \(q^{m-1}\).  Its
parenthesized factor tends to one, so the finite-ramification normal form of
\cref{thm:uq-finite-ramification} gives exactly \(m-1\) inverse
determinations.
\end{proof}

\section{\(q\)-factorials and their exact ramification}
\label{sec:qs-qfactorials}

For \(n\in\NaturalNumbers\), put
\[
  [n]_q!:=\prod_{j=1}^{n}[j]_q.
\]

\begin{proposition}[Zero-base factorial inverse]
\label{prop:qs-qfactorial-zero}
For \(n\ge2\),
\[
  [n]_q!
  =1+(n-1)q+\frac{(n-2)(n+1)}2q^2+\BigO(q^3).
\]
Consequently, with \(u=y-1\),
\[
  q=\frac{u}{n-1}
   -\frac{(n-2)(n+1)}{2(n-1)^3}u^2+\BigO(u^3).
\]
\end{proposition}

\begin{proof}
In the product
\[
  [n]_q!=\prod_{j=2}^{n}(1+q+\cdots+q^{j-1}),
\]
the coefficient of \(q\) is \(n-1\).  A term of degree two is obtained
either by choosing \(q\) from two distinct factors, in
\(\binom{n-1}{2}\) ways, or by choosing \(q^2\) from one of the
\(n-2\) factors with \(j\ge3\).  Their sum is
\((n-2)(n+1)/2\).  Reversion of
\(u=(n-1)q+Aq^2+\cdots\) proves the inverse formula.
\end{proof}

\begin{theorem}[Exact root-of-unity ramification of \(q\)-factorials]
\label{thm:qs-qfactorial-root}
Let \(\zeta\ne1\) be a primitive \(d\)-th root of unity.  Then
\[
  \operatorname{ord}_{q=\zeta}[n]_q!
  =\left\lfloor\frac nd\right\rfloor.
\]
In particular, if \(M=\lfloor n/2\rfloor\),
\[
  [n]_q!=M!(q+1)^M\bigl(1+\BigO(q+1)\bigr)
  \qquad(q\to-1),
\]
and the inverse at target zero has the \(M\) branches
\[
  q=-1+\omega_M^j\left(\frac{y}{M!}\right)^{1/M}
    +\BigO(y^{2/M}),
  \qquad 0\le j<M.
\]
\end{theorem}

\begin{proof}
Because \(\zeta\ne1\), the denominator in
\([j]_q=(1-q^j)/(1-q)\) is nonzero at \(\zeta\).  The numerator has a
simple zero there exactly when \(d\mid j\).  Exactly
\(\lfloor n/d\rfloor\) factors therefore vanish, each simply.

For \(d=2\), the even factors satisfy
\[
  [2r]_q=r(q+1)+\BigO((q+1)^2),
\]
while the odd factors tend to one.  Multiplying \(r=1,\ldots,M\) gives the
coefficient \(M!\).  Applying
\cref{thm:uq-finite-ramification} to the resulting exact zero order gives the
Puiseux branches and their remainder.
\end{proof}

Writing \(N=\binom n2\), one also has
\[
  [n]_q!=q^N\left(1+\frac{n-1}{q}+\BigO(q^{-2})\right),
\]
and the same logarithmic coefficient comparison as in
\cref{thm:gq-large-inverse} yields, on each root branch,
\[
  q=y^{1/N}-\frac{n-1}{N}+\BigO(y^{-1/N})
   =y^{1/N}-\frac2n+\BigO(y^{-1/N}).
\]

\section{The \(q\)-gamma function}
\label{sec:qs-qgamma}

For \(0<q<1\) and \(x>0\), define
\[
  \Gamma_q(x)
  =(1-q)^{1-x}\frac{(q;q)_\infty}{(q^x;q)_\infty}.
\]
It satisfies
\[
  \Gamma_q(x+1)=[x]_q\Gamma_q(x),
  \qquad
  \Gamma_q(1)=\Gamma_q(2)=1.
\]

\begin{proposition}[Exact generalized expansion at \(q=0\)]
\label{prop:qs-qgamma-zero}
For \(0<q<1\) and \(x>0\),
\[
  \log\Gamma_q(x)
  =\sum_{m\ge1}\frac{q^{mx}}{m(1-q^m)}
   -\sum_{m\ge1}\frac{q^m}{m(1-q^m)}
   -(1-x)\sum_{m\ge1}\frac{q^m}{m}.
\]
All three series converge absolutely.  Expanding
\((1-q^m)^{-1}=\sum_{r\ge0}q^{mr}\) gives a convergent generalized power
series with exponents in
\[
  \{m(x+r):m\ge1,\ r\ge0\}\cup\NaturalNumbers.
\]
For rational \(x\) it is a Puiseux series; for irrational \(x\), its natural
ordering is a generalized-power ordering rather than an ordinary Taylor
ordering.  The two primitive contributions are
\[
  \log\Gamma_q(x)=q^x-(2-x)q+\text{higher supported powers}.
\]
\end{proposition}

\begin{proof}
Take logarithms in the product definition and use the absolutely convergent
Lambert identity
\[
  \log(a;q)_\infty
  =-\sum_{m\ge1}\frac{a^m}{m(1-q^m)}
\]
from \cref{lem:ip-lambert-log}, together with
\(\log(1-q)=-\sum_{m\ge1}q^m/m\).  Absolute convergence permits expansion of
each geometric denominator and rearrangement.  The displayed exponent set is
then immediate, and the terms \(m=1,r=0\) and the integer exponent one give
\(q^x-(2-x)q\), subject to the stated collisions at special \(x\).
\end{proof}

The logarithmic derivative and its derivative are
\[
  \psi_q(x)
  =-\log(1-q)+(\log q)
    \sum_{j\ge0}\frac{q^{x+j}}{1-q^{x+j}},
\]
\[
  \psi_q'(x)
  =(\log q)^2
    \sum_{j\ge0}\frac{q^{x+j}}{(1-q^{x+j})^2}>0.
\]
Termwise differentiation is justified locally uniformly for \(x>0\).

\begin{theorem}[Unique \(q\)-gamma minimum and two real inverse branches]
\label{thm:qs-qgamma-minimum}
For every \(0<q<1\), there is a unique \(x_q>0\) with
\(\psi_q(x_q)=0\).  The function \(\Gamma_q\) strictly decreases on
\((0,x_q)\), strictly increases on \((x_q,\infty)\), and tends to infinity
at both endpoints.  Thus each target strictly above
\(\Gamma_q(x_q)\) has exactly two positive argument inverses, the minimum has
one double inverse point, and smaller targets have none.  The two branches
meet in a square-root fold.
\end{theorem}

\begin{proof}
The formula above makes \(\psi_q\) strictly increasing.  As \(x\downarrow0\),
the \(j=0\) summand is asymptotic to
\[
  (\log q)\frac1{-x\log q}=-\frac1x,
\]
so \(\psi_q(x)\to-\infty\).  As \(x\to\infty\), the series tends to zero and
\(\psi_q(x)\to-\log(1-q)>0\).  Hence it has one zero.

Since \(\Gamma_q'/\Gamma_q=\psi_q\), the derivative has the asserted signs.
The factor \((q^x;q)_\infty\) has a simple zero at \(x=0\), so
\(\Gamma_q(x)\to\infty\) there.  As \(x\to\infty\), the denominator tends to
one while \((1-q)^{1-x}\to\infty\).  The branch count follows by continuity
and strict monotonicity.  At the minimum,
\[
  (\log\Gamma_q)''(x_q)=\psi_q'(x_q)>0,
\]
so Taylor's theorem and \cref{thm:uq-finite-ramification} give the
square-root fold.
\end{proof}

\begin{theorem}[All-order \(q\to1\) expansion of \(\Gamma_q\)]
\label{thm:qs-qgamma-one}
Let \(q=\EulerE^{-t}\), with \(t\to0\) in a closed subsector of
\(|\arg t|<\pi/2\).  Uniformly for \(x\) in compact subsets of
\(\Re x>0\),
\[
  \log\Gamma_{\EulerE^{-t}}(x)
  \sim\log\Gamma(x)+c_1(x)t+\sum_{r\ge1}c_{2r}(x)t^{2r},
\]
where
\[
  c_1(x)=\frac{(x-1)(2-x)}4
\]
and
\[
  c_{2r}(x)
  =\frac{B_{2r}}{2r(2r)!}
   \left(1-x+\frac{B_{2r+1}(x)}{2r+1}\right).
\]
No odd algebraic power beyond \(t^1\) occurs.
\end{theorem}

\begin{proof}
Insert the all-order coalescing expansion of
\((q^x;q)_\infty\) from \cref{thm:ip-qx-expansion} into
\[
  \log\Gamma_q(x)
  =(1-x)\log(1-q)+\log(q;q)_\infty-\log(q^x;q)_\infty.
\]
Use the \(x=1\) specialization for \((q;q)_\infty\) and
\[
  \log\frac{1-\EulerE^{-t}}{t}
  =-\frac t2+\sum_{r\ge1}
   \frac{B_{2r}}{2r(2r)!}t^{2r}.
\]
The \(t^{-1}\), logarithmic, and constant normalization terms cancel,
leaving \(\log\Gamma(x)\).  The linear coefficient simplifies to
\[
  \frac1{24}-\frac{B_2(x)}4-\frac{1-x}{2}
  =\frac{(x-1)(2-x)}4.
\]
At every even order, the \(x=1\) Bernoulli-polynomial term vanishes and the
remaining two terms combine to \(c_{2r}(x)\).  The expansions invoked above
are locally uniform with uniform differentiated remainders, so subtraction
preserves the asserted uniform asymptotic meaning.
\end{proof}

In particular,
\[
  c_2(x)=\frac{(x-2)(x-1)(2x+3)}{144},
\]
and, for fixed \(x\notin\{1,2\}\), if
\[
  v=\log\frac{y}{\Gamma(x)},
\]
ordinary reversion gives
\[
  t=\frac{v}{c_1(x)}
   -\frac{c_2(x)}{c_1(x)^3}v^2
   +\frac{2c_2(x)^2}{c_1(x)^5}v^3+\BigO(v^4),
  \qquad q=\EulerE^{-t}.
\]
At \(x=1,2\), \(\Gamma_q(x)\equiv1\), so there is no base inverse.

\section{\(q\)-digamma inversion}
\label{sec:qs-qdigamma}

\begin{proposition}[Global argument inverse and the exceptional base fold]
\label{prop:qs-qdigamma-inverse}
For fixed \(0<q<1\), the function \(\psi_q\) is a strictly increasing
bijection
\[
  (0,\infty)\longrightarrow(-\infty,-\log(1-q)).
\]
For \(q=\EulerE^{-t}\),
\[
  \psi_{\EulerE^{-t}}(x)
  \sim\psi(x)+\frac{3-2x}{4}t
  +\sum_{r\ge1}c_{2r}'(x)t^{2r}.
\]
At fixed \(x\ne3/2\), base inversion at \(q=1\) is ordinary.  At \(x=3/2\),
\[
  \psi_{\EulerE^{-t}}(3/2)-\psi(3/2)
  =-\frac{t^2}{288}+\BigO(t^4),
\]
so the base inverse is a quadratic fold.
\end{proposition}

\begin{proof}
Strict increase and the two endpoint limits were proved in
\cref{thm:qs-qgamma-minimum}; these give the bijection.  Differentiate
\cref{thm:qs-qgamma-one} termwise in \(x\), which is permitted by its compact
local uniformity.  The linear coefficient
\((3-2x)/4\) vanishes only at \(3/2\).  Direct differentiation of
\[
  c_2(x)=\frac{(x-2)(x-1)(2x+3)}{144}
\]
gives \(c_2'(3/2)=-1/288\).  The inverse function theorem handles the
nonexceptional points, and the finite-ramification theorem handles the
nondegenerate quadratic point.
\end{proof}

\section{\(q\)-beta coordinate and base inverses}
\label{sec:qs-qbeta}

For \(x,y>0\), define
\[
  B_q(x,y):=\frac{\Gamma_q(x)\Gamma_q(y)}{\Gamma_q(x+y)}.
\]

\begin{theorem}[Coordinate monotonicity and base type change]
\label{thm:qs-qbeta}
For every \(0<q<1\), \(B_q(x,y)\) is strictly decreasing in either positive
coordinate while the other is fixed.  Hence each coordinate inverse is
unique on its range.

As \(q=\EulerE^{-t}\to1\),
\[
  \log\frac{B_{\EulerE^{-t}}(x,y)}{B(x,y)}
  =b_1(x,y)t+b_2(x,y)t^2+\BigO(t^4),
\]
where
\[
  b_1(x,y)=\frac{xy-1}{2},
  \qquad
  b_2(x,y)=-\frac{x^2y+xy^2-xy-1}{24}.
\]
The base inverse is ordinary when \(xy\ne1\).  On \(xy=1\),
\[
  b_2(x,1/x)=-\frac{(x-1)^2}{24x};
\]
there is a quadratic fold for \(x\ne1\), while
\(B_q(1,1)\equiv1\) makes the base completely unidentifiable at
\((1,1)\).
\end{theorem}

\begin{proof}
The logarithmic coordinate derivatives are
\[
  \partial_x\log B_q(x,y)=\psi_q(x)-\psi_q(x+y)<0,
\]
\[
  \partial_y\log B_q(x,y)=\psi_q(y)-\psi_q(x+y)<0,
\]
because \(\psi_q\) is strictly increasing.  This proves strict coordinate
monotonicity and uniqueness of each coordinate inverse.

Subtract the \(x+y\) instance of \cref{thm:qs-qgamma-one} from the sum of its
\(x\) and \(y\) instances.  Elementary expansion of \(c_1\) and \(c_2\)
gives the displayed \(b_1,b_2\); no cubic term exists because the
\(q\)-gamma expansion has no odd term beyond the linear one.  The inverse
function theorem applies when \(b_1\ne0\).  On \(xy=1\), substitution gives
the stated negative \(b_2\), nonzero unless \(x=1\).  At the remaining point,
\(\Gamma_q(1)=\Gamma_q(2)=1\) proves the exact identity.
\end{proof}

\begin{warning}
The claim that \(q\)-beta is always strictly decreasing in its base is false:
\(B_q(1,1)=1\) for every base.  The theorem above is the corrected statement.
It proves strictness in the two shape coordinates and gives the exact
base-inverse type change.
\end{warning}

\section{Euler \(q\)-exponentials}
\label{sec:qs-qexponentials}

It is useful to separate the product variable from the classical-limit
scaling:
\[
  \mathsf e_q(z):=\frac1{(z;q)_\infty},
  \qquad
  \mathsf E_q(z):=(-z;q)_\infty,
\]
\[
  e_q(x):=\mathsf e_q((1-q)x),
  \qquad
  E_q(x):=\mathsf E_q((1-q)x).
\]

\begin{proposition}[Exact inverse reduction and \(q\to1\) generator]
\label{prop:qs-qexponential}
If \(A_q\) is a selected argument-inverse branch of
\(a\mapsto(a;q)_\infty\), then
\[
  \mathsf e_q^{-1}(y)=A_q(y^{-1}),
  \qquad
  \mathsf E_q^{-1}(y)=-A_q(y).
\]
Writing \(h=1-q\), the scaled functions satisfy the convergent logarithmic
identities
\[
  \log e_q(x)
  =\sum_{m\ge1}\frac{h^m x^m}{m(1-(1-h)^m)},
\]
\[
  \log E_q(x)
  =\sum_{m\ge1}\frac{(-1)^{m+1}h^m x^m}
                       {m(1-(1-h)^m)}.
\]
At each fixed order in \(h\), only finitely many indices \(m\) contribute.
\end{proposition}

\begin{proof}
The inverse formulas are immediate after rewriting the two target equations
as \((z;q)_\infty=y^{-1}\) and \((-z;q)_\infty=y\), with the selected branch
data retained.  Apply \cref{lem:ip-lambert-log} with
\(a=\pm hx\) and \(q=1-h\), remembering the minus sign introduced by the
reciprocal in \(e_q\).  Since the \(m\)-th summand begins at order
\(h^{m-1}\) after the denominator is expanded, only finitely many \(m\) can
affect any prescribed coefficient.
\end{proof}

For example,
\[
  \log e_q(x)
  =x+\frac{x^2}{4}h+
   \left(\frac{x^2}{8}+\frac{x^3}{9}\right)h^2+\BigO(h^3),
\]
\[
  \log E_q(x)
  =x-\frac{x^2}{4}h+
   \left(-\frac{x^2}{8}+\frac{x^3}{9}\right)h^2+\BigO(h^3).
\]
Bell exponentiation gives the functions themselves, while
\(\mathsf e_q(z)\mathsf E_q(-z)=1\) is exact.

\section{Basic hypergeometric argument inversion}
\label{sec:qs-hypergeometric}

\begin{theorem}[Universal local series-variable inverse]
\label{thm:qs-hypergeometric-inverse}
Let a basic hypergeometric function, after all removable parameter
cancellations, have a convergent local series
\[
  F(z;\bm a,\bm b,q)
  =1+c_1z+c_2z^2+c_3z^3+\cdots,
  \qquad c_1\ne0.
\]
Then its argument inverse at \(y=1\) is
\[
  z=\frac{y-1}{c_1}
   -\frac{c_2}{c_1^3}(y-1)^2
   +\left(\frac{2c_2^2}{c_1^5}-\frac{c_3}{c_1^4}\right)(y-1)^3
   +\BigO((y-1)^4).
\]
All higher coefficients are given by the analytic Lagrange--Bürmann formula.
\end{theorem}

\begin{proof}
The condition \(c_1\ne0\) gives a local analytic inverse.  Put
\(u=y-1\) and
\[
  z=A u+B u^2+C u^3+\BigO(u^4).
\]
Substitution in \(u=c_1z+c_2z^2+c_3z^3+\cdots\) yields
\[
  c_1A=1,\qquad
  c_1B+c_2A^2=0,\qquad
  c_1C+2c_2AB+c_3A^3=0.
\]
Solving gives the displayed coefficients.  The all-order statement is
\cref{thm:uq-analytic-lagrange}.
\end{proof}

If a numerator parameter reaches \(q^{-N}\), the series truncates and the
argument inverse becomes an algebraic covering.  If a denominator parameter
approaches the same truncation locus, common finite \(q\)-Pochhammer factors
must be cancelled before applying the theorem; otherwise a removable
singularity is misclassified as a branch point.

\section{Multiplicative theta products}
\label{sec:qs-theta}

For \(0<|q|<1\), let
\[
  \vartheta(z;q)
  :=(z;q)_\infty(q/z;q)_\infty(q;q)_\infty,
  \qquad z\in\ComplexNumbers^\times.
\]

\begin{proposition}[Simple theta zero atlas]
\label{prop:qs-theta-zeros}
The zeros of \(z\mapsto\vartheta(z;q)\) are exactly
\(z=q^m\), \(m\in\IntegerNumbers\), and every zero is simple.  Hence the inverse at target
zero is ordinary on every sheet:
\[
  z=q^m+\frac{y}{\partial_z\vartheta(q^m;q)}+\BigO(y^2).
\]
Moreover,
\[
  \vartheta(qz;q)=-z^{-1}\vartheta(z;q),
\]
so one zero derivative determines all the others.
\end{proposition}

\begin{proof}
The factor \((z;q)_\infty\) has one simple vanishing factor precisely at
\(z=q^{-j}\), \(j\ge0\).  The factor \((q/z;q)_\infty\) has one simple
vanishing factor precisely at \(z=q^{j+1}\), \(j\ge0\).  Because
\(|q|<1\), the two lists do not overlap, and every other factor is nonzero.
Their union is \(\{q^m:m\in\IntegerNumbers\}\), proving the zero and multiplicity claims.
The local inverse formula is the analytic inverse function theorem.

Finally,
\[
  (qz;q)_\infty=\frac{(z;q)_\infty}{1-z},
  \qquad
  (1/z;q)_\infty=(1-z^{-1})(q/z;q)_\infty.
\]
Multiplying and using
\((1-z^{-1})/(1-z)=-z^{-1}\) proves quasi-periodicity.  Differentiating it
at a zero transports the derivative normalization.
\end{proof}
